% 03_spo_criteria.tex — ตัวบ่งชี้ 1.1 เกณฑ์ สป.อว. + เปิด AUN-QA Criteria
\section{ตัวบ่งชี้ที่ 1.1 การบริหารจัดการหลักสูตรตามเกณฑ์มาตรฐาน}

หลักสูตรวิศวกรรมศาสตรมหาบัณฑิต สาขาวิชาวิศวกรรมคอมพิวเตอร์และปัญญาประดิษฐ์
ดำเนินการตามเกณฑ์มาตรฐานหลักสูตรระดับบัณฑิตศึกษา พ.ศ. 2558 และเกณฑ์ AUN-QA Version 4.0
โดยมีการพัฒนาหลักสูตรอย่างเป็นระบบตั้งแต่การสำรวจความต้องการของผู้มีส่วนได้ส่วนเสีย
การกำหนดผลการเรียนรู้ที่คาดหวัง การออกแบบรายวิชา จนถึงการประเมินคุณภาพหลักสูตร

\begin{tabularx}{\linewidth}{|c|>{\raggedright\arraybackslash}X|c|}
    \hline
    \rowcolor{gray!15}
    \textbf{Criteria} & \textbf{หัวข้อการประเมิน} & \textbf{คะแนน (1--7)} \\
    \hline
    1 & Expected Learning Outcomes & 3 \\
    \hline
    2 & Programme Structure and Content & 3 \\
    \hline
    3 & Teaching and Learning Approach & 2 \\
    \hline
    4 & Student Assessment & 3 \\
    \hline
    5 & Academic Staff Quality & 2 \\
    \hline
    6 & Student Support Services & 2 \\
    \hline
    7 & Facilities and Infrastructure & 3 \\
    \hline
    8 & Output, Outcomes and Impact & 1 \\
    \hline
    \textbf{รวมเฉลี่ย} & & \textbf{2.4} \\
    \hline
\end{tabularx}

\marknote{คะแนนข้างต้นเป็นการประเมินตนเองเบื้องต้น หลักสูตรใหม่ที่ยังไม่มีบัณฑิตจะได้คะแนนสูงสุดไม่เกิน 3/7 ในเกณฑ์ที่ต้องการผลลัพธ์จริง}
